\documentclass[11pt]{article}
\usepackage[a4paper,margin=0.87in]{geometry}
\usepackage[T1]{fontenc}
\usepackage{lmodern,amsmath,amssymb,amsthm,mathtools,microtype,booktabs}
\usepackage{enumitem,fancyhdr,xcolor,hyperref}
\hypersetup{colorlinks=true,linkcolor=black,citecolor=black,urlcolor=blue,
 pdftitle={A four-point counterexample to the verbose pullback triangle inequality},
 pdfauthor={Siavash Sadeghi (contributor and submitter); AI-assisted research note},pdfsubject={Proposed resolution of original AIM 544, current AIM 534}}
\setlength{\parindent}{0pt}
\setlength{\parskip}{5pt plus 1pt}
\setlength{\headheight}{14pt}
\pagestyle{fancy}
\fancyhf{}
\lhead{\small AIM 544 (now 534) \;|\; Proposed counterexample}
\rhead{\small 2 October 2026}
\cfoot{\thepage}
\newtheorem{theorem}{Theorem}
\newtheorem{lemma}[theorem]{Lemma}
\newtheorem{proposition}[theorem]{Proposition}
\newtheorem{corollary}[theorem]{Corollary}
\theoremstyle{remark}
\newtheorem{remark}[theorem]{Remark}
\newcommand{\F}{\mathbb F}
\newcommand{\VR}{\operatorname{VR}}
\newcommand{\rk}{\operatorname{rank}}
\newcommand{\diam}{\operatorname{diam}}
\newcommand{\sq}{\sqcup}
\newcommand{\barc}[2]{\{(#1,#1)\}^{#2}}
\title{\vspace{-1.1em}A four-point counterexample to the\newline verbose pullback triangle inequality}
\author{Siavash Sadeghi\\\small Contributor and submitter; AI-assisted research note}
\date{2 October 2026}
\begin{document}
\maketitle
\vspace{-1.2em}
\begin{abstract}
The fixed-degree pullback matching distance of verbose Vietoris--Rips barcodes
fails the triangle inequality. In degree one, three four-point ultrametric
spaces, with nonzero distances only $1$ and $2$, have pairwise distances
$0,0,1$. More generally, in every degree $k\ge1$ and over every field, the
same phenomenon occurs on three spaces of equal cardinality $k+3$.
Every barcode involved is nonempty. A rank-counting argument proves the
positive lower bound over \emph{all} finite common pullbacks, not merely
an enumerated set. The common cardinality $k+3$ is minimal. The degree-one
zero-distance pair used here is already in M\'emoli--Zhou, Example~5.10;
the third space and the uniform lower bound complete a counterexample.
\end{abstract}

\textbf{Evidence status.} Complete proposed disproof, awaiting independent
mathematical review. The supplied proof and programs are AI-generated. Codex and separate
AI reviewers checked the argument, programs, and target sources during
submission preparation; details appear in \texttt{verification\_record.md}.
No human peer review, formal proof-assistant verification, or repository
acceptance is claimed.

\section{The target and conventions}
The target is original AIM 544, \emph{Triangle inequality for the pullback
distance of verbose persistence barcodes}, pinned at \texttt{aa776a0}~\cite{aim}.
As of 2 October 2026 it is AIM 534 at \texttt{fa98b75}; only the numbered
heading changed, and the target remains Open. Full revision hashes and the
dated prior-submission audit are in \texttt{target\_audit.md}. No overlapping
claim was found in the inspected public records.

Fix a field $\F$. For a finite nonempty pseudometric space $(S,d)$, give each
simplex its diameter as filtration value. Write $V_k(S,d)$ for the degree-$k$
barcode obtained by the ordinary boundary reduction \emph{retaining all
zero-length pairs}. Endpoint pairs have their full multiplicities.
For equal-size finite barcodes, $m$ is the least maximum $\ell^\infty$
endpoint discrepancy over bijections; no additional diagonal points are
allowed. For finite metric spaces $A,B$, the target distance is
\begin{equation}
D_k(A,B)=\inf_{\substack{S\text{ finite}\,;\ p:S\twoheadrightarrow A\,;\ q:S\twoheadrightarrow B}}
 m\bigl(V_k(S,p^*d_A),V_k(S,q^*d_B)\bigr).
\label{eq:target}
\end{equation}
Distinct vertices at zero pullback distance are not identified. We work
only in degrees $k\ge1$, so all bars have finite endpoints. On a common
set of $s$ vertices, either pullback has exactly $\binom{s-1}{k+1}$
degree-$k$ pairs, since its final complex is the full simplex. Thus the
compared barcodes have equal cardinality (Lemma~\ref{lem:rank}). The published
fixed-degree question is in Remark~1.1 and Definition~5.3 of~\cite{mz}.

\section{The four-point counterexample}
Take three copies of $\{1,2,3,4\}$ with distance matrices
\begin{equation}
 d_X=\begin{pmatrix}
0&1&1&2\\1&0&1&2\\1&1&0&2\\2&2&2&0
\end{pmatrix},\quad
 d_Y=\begin{pmatrix}
0&1&2&2\\1&0&2&2\\2&2&0&1\\2&2&1&0
\end{pmatrix},\quad
 d_Z=\begin{pmatrix}
0&2&2&2\\2&0&2&2\\2&2&0&2\\2&2&2&0
\end{pmatrix}.
\label{eq:matrices}
\end{equation}
All are genuine ultrametrics. At scale $1$, the clusters of $X$ have sizes
$3,1$, those of $Y$ have sizes $2,2$, and those of $Z$ have sizes $1,1,1,1$.
At scale $2$, all four vertices form a simplex.

\begin{theorem}\label{thm:four}
For the spaces in~\eqref{eq:matrices}, over every field,
\[
\boxed{D_1(X,Y)=0,\qquad D_1(Y,Z)=0,\qquad D_1(X,Z)=1.}
\]
In particular, these spaces disprove the universal assertion in AIM 544.
\end{theorem}

We first give the elementary counting facts that justify every barcode
and the lower bound over arbitrary pullbacks.

\begin{lemma}[Boundary ranks of a simplex]\label{lem:rank}
For the full simplex on $r\ge1$ vertices and $k\ge1$,
\[
\rk\partial_{k+1}=\binom{r-1}{k+1},
\]
where a binomial coefficient is zero when its lower index exceeds its
nonnegative upper index. For a disjoint union of simplexes on
$r_1,\ldots,r_s$ vertices, the rank is the sum of these quantities.
\end{lemma}
\begin{proof}
Fix one vertex $v$. The boundaries of the $(k+1)$-simplexes containing
$v$ are independent: each has a distinct face of dimension $k$ not
containing $v$, with coefficient $1$ or $-1$. They span the boundary
space. Indeed, the boundary of a $(k+1)$-simplex not containing $v$ is a
linear combination of these boundaries, obtained by applying
$\partial^2=0$ to its cone with vertex $v$. There are
$\binom{r-1}{k+1}$ such simplexes. The argument works in every characteristic.
Disjoint components give a direct sum.
\end{proof}

\begin{lemma}[Counting the diagonal bars]\label{lem:counts}
Suppose a finite simplicial filtration changes only at
$t_0<\cdots<t_s$, and its positive-degree homology vanishes at every
threshold. Then every positive-degree verbose bar is diagonal.
The number of degree-$k$ bars with death time at most $t_i$ is
$\rk\partial_{k+1}(K_{t_i})$.
Consequently, the multiplicity of $(t_i,t_i)$ is the increment of this
rank from the preceding threshold.
\end{lemma}
\begin{proof}
A positive-length bar would give nonzero homology at some threshold or
on an interval between successive thresholds. Thus no such bars occur.
In a boundary reduction, nonzero columns of degree $k+1$ are exactly the
deaths of degree-$k$ pairs. Reducing all columns with entry time at most
$t_i$ gives as many nonzero columns as the rank of the corresponding
boundary map. This statement includes deaths paired with births at the
same filtration value. Differences count deaths at the individual
thresholds.
\end{proof}

\subsection{The two zero distances}
The filtrations here are disjoint unions of simplexes at every scale, so
Lemmas~\ref{lem:rank}--\ref{lem:counts} give
\begin{equation}
V_1(X)=\{(1,1),(2,2),(2,2)\},\qquad
V_1(Y)=V_1(Z)=\barc{2}{3}.
\label{eq:base}
\end{equation}
Use the original four-element set, with bijections onto $Y$ and $Z$.
Their equal barcodes prove $D_1(Y,Z)=0$.

For $X,Y$, use a five-element set $S=\{1,2,3,4,5\}$ and the surjections
\begin{equation}
\begin{array}{c|ccccc}
s&1&2&3&4&5\\\hline
p(s)\in X&1&2&3&4&4\\
q(s)\in Y&1&2&3&4&1
\end{array}.
\label{eq:witness}
\end{equation}
The first pullback repeats the isolated point of $X$; the second repeats
one member of a close pair in $Y$. Their cluster sizes at the thresholds
are
\[
\begin{array}{c|ccc}
&0&1&2\\\hline
(S,p^*d_X)&2,1,1,1&3,2&5\\
(S,q^*d_Y)&2,1,1,1&3,2&5
\end{array}.
\]
For degree one, the boundary ranks at these thresholds are $0,1,6$.
Hence both pullbacks have barcode
\begin{equation}
\{(1,1)\}\sq\barc{2}{5}.
\label{eq:pullback}
\end{equation}
This is a zero-cost admissible matching, proving $D_1(X,Y)=0$.

\subsection{A lower bound for every common pullback}
Let $S$ be \emph{any} finite set surjecting onto $X$ and $Z$.
Let $a,b,c,d\ge1$ be the fiber sizes for its map to $X$.
At threshold $0$, the $X$-pullback is the disjoint union of simplexes
on $a,b,c,d$ vertices. At threshold $1$, these merge into simplexes
on $a+b+c$ and $d$ vertices. The number of $(1,1)$ bars is therefore
\begin{align}
&\binom{a+b+c-1}{2}+\binom{d-1}{2}
 -\binom{a-1}{2}-\binom{b-1}{2}-\binom{c-1}{2}-\binom{d-1}{2}\notag\\
&\hspace{2cm}=ab+ac+bc-2\ \ge\ 1.
\label{eq:positive}
\end{align}
This includes all possible duplication patterns, without any size bound.

The $Z$-pullback has only distances $0$ and $2$. At threshold $0$ it is a
disjoint union of simplexes, and at threshold $2$ it is a full simplex.
Its degree-one bars are therefore only $(0,0)$ and $(2,2)$.
Every bijection must match the $(1,1)$ bar supplied by~\eqref{eq:positive}
to one of these two pairs. Either choice has cost $1$.
Thus every admissible comparison in~\eqref{eq:target} costs at least $1$.

Conversely, the original four-point barcodes in~\eqref{eq:base} can be
matched with cost $1$. This proves $D_1(X,Z)=1$ and completes the proof of
Theorem~\ref{thm:four}.
\qed

\paragraph{Attribution and scope.}
The spaces $X,Y$, their degree-one barcodes, and a zero-cost duplication
comparison already occur in M\'emoli--Zhou, Example~5.10 and Figure~7,
after relabeling~\cite{mz}. No novelty is claimed for that pair. The choice
of an equilateral third space with distance $2$, together with the
uniform obstruction~\eqref{eq:positive}, completes the proposed disproof.
Our $Z$ differs from the equilateral distance-$1$ space denoted $Z$ in
that source example.
This is a fixed-degree argument, not the source's distinct all-degree
pullback-interleaving counterexample. It does not depend on any convention
for empty barcodes: every original barcode here has three entries and
every common pullback has at least four vertices.

\section{A counterexample in every positive degree}
\begin{theorem}\label{thm:general}
For every field $\F$ and every integer $k\ge1$, there are three ultrametric
spaces $X_k,Y_k,Z_k$, each with $k+3$ points and nonzero distances in
$\{1,2\}$, such that
\[
D_k(X_k,Y_k)=D_k(Y_k,Z_k)=0,\qquad D_k(X_k,Z_k)=1.
\]
The degree-$k$ barcode of every original space and every common pullback
in these comparisons is nonempty.
\end{theorem}
\begin{proof}
Set $n=k+3$. Give $X_k$ a close cluster of size $k+2$ and a singleton,
and give $Y_k$ two close clusters of sizes $k+1$ and $2$.
Distances between distinct points in the same close cluster are $1$;
all remaining off-diagonal distances are $2$.
Let $Z_k$ be the equilateral $n$-point space with distance $2$.

Lemmas~\ref{lem:rank}--\ref{lem:counts} give
\begin{align}
V_k(X_k)&=\{(1,1)\}\sq\barc{2}{k+1},\label{eq:general-base}\\
V_k(Y_k)&=V_k(Z_k)=\barc{2}{k+2}.
\end{align}
In particular, the original sets give a zero-cost comparison of $Y_k,Z_k$.

Repeat the singleton in $X_k$ and one vertex of the size-$(k+1)$ cluster
in $Y_k$. Both pullbacks have $k+4$ vertices. At $0$, one fiber has size
$2$ and all others have size $1$; since $k\ge1$, this yields no degree-$k$
bars at $0$. At $1$, both have component sizes $k+2,2$, with boundary
rank $1$. At $2$, the rank is $\binom{k+3}{k+1}$. Both barcodes are
\begin{equation}
\{(1,1)\}\sq\barc{2}{\binom{k+3}{k+1}-1},
\end{equation}
so $D_k(X_k,Y_k)=0$.

Now consider any pullback of $X_k$. Let $m_1,\ldots,m_{k+2}\ge1$
be the fiber sizes over its large cluster, and put $M=\sum_i m_i$.
The singleton fiber's contribution cancels between levels $0$ and $1$.
The multiplicity of $(1,1)$ is
\begin{equation}
B=\binom{M-1}{k+1}-\sum_{i=1}^{k+2}\binom{m_i-1}{k+1}\ \ge\ 1.
\label{eq:all-multiplicities}
\end{equation}
For a direct proof of the inequality, form disjoint sets $A_i$ with
$|A_i|=m_i-1$ and a further disjoint set $R$ of size $k+1$.
Their union has size $M-1$. The first binomial counts all its
$(k+1)$-element subsets; the sum subtracts just those lying entirely
in some $A_i$. The subset $R$ is not subtracted, so the remainder is
at least one.

Every pullback of $Z_k$ has only $(0,0)$ and $(2,2)$ bars, as in the
degree-one argument. Equation~\eqref{eq:all-multiplicities} forces matching
cost at least $1$ for every common pullback. The original barcodes
in~\eqref{eq:general-base} give cost $1$, proving equality.
Finally, on any $s\ge k+3$ vertices a full simplex has
$\binom{s-1}{k+1}>0$ degree-$k$ pairs, which proves nonemptiness.
\end{proof}

\section{Minimal equal cardinality and consequences}
\begin{lemma}[A diameter bar]\label{lem:diameter}
Let $A$ be a finite pseudometric space with $r\ge k+2$ distinct vertices
and diameter $\delta$, where $k\ge1$. Then $V_k(A)$ contains
$(\delta,\delta)$.
\end{lemma}
\begin{proof}
Choose two vertices realizing the diameter and extend them to a set of
$k+2$ vertices, defining a $(k+1)$-simplex $\sigma$.
The cycle $\partial\sigma$ has a nonzero coefficient on a $k$-face
containing those two vertices: delete any of the other $k$ vertices.
That face has entry time $\delta$. Thus $\partial\sigma$ is not in the
chain space strictly below $\delta$, and the dimension of the space of
$k$-cycles increases at $\delta$.

The number of degree-$k$ barcode births through a threshold is the
number of zero reduced degree-$k$ columns, namely the dimension of
that threshold's $k$-cycle space. Hence at least one bar is born at
$\delta$. The filtration is the full simplex at $\delta$, so every
bar born then also dies then. If $\delta=0$, the same argument uses
the empty complex below $0$.
\end{proof}

\begin{proposition}\label{prop:minimal}
Among triples whose three spaces have the same cardinality, $k+3$ is the
smallest cardinality at which the triangle inequality for $D_k$ can fail.
\end{proposition}
\begin{proof}
For equal cardinality $n\le k+1$, the original degree-$k$ barcodes are
empty, so $D_k$ is identically zero on that cardinality class, using the
repository's explicit empty-matching convention.

For $n=k+2$, the original degree-$k$ barcode has exactly one pair by
Lemma~\ref{lem:rank}. Lemma~\ref{lem:diameter} identifies it as
$(\diam A,\diam A)$. This gives
\[
D_k(A,B)\le |\diam A-\diam B|.
\]
Every admissible common pullback has at least $k+2$ vertices, has the
same diameter as its target, and has a diameter bar by
Lemma~\ref{lem:diameter}. All endpoints in the other barcode are at most
its target's diameter. Matching the larger diameter bar therefore costs
at least $|\diam A-\diam B|$, proving equality. The absolute-value triangle
inequality applies on this cardinality class. Theorem~\ref{thm:general}
attains failure at $n=k+3$.
\end{proof}

\paragraph{Consequences.}
In every positive degree, the zero-distance relation is not transitive.
Moreover, there is no finite constant $C$ such that
\[
D_k(A,C')\le C\bigl(D_k(A,B)+D_k(B,C')\bigr)
\]
for all finite metric spaces: the right-hand side for our triple is zero.
Multiplying all distances by $a>0$ gives distances $0,0,a$, so there is
no universal additive bound on the triangle defect either.
Taking an infimum of sums along chains produces a \emph{different}
quantity, which identifies $X_k$ and $Z_k$ at zero cost.

The witness also works if common pullbacks are restricted to
\emph{correspondences}, meaning $s\mapsto(p(s),q(s))$ is injective.
In the general construction, take identity labels for the first $n$
vertices and an additional pair $(n,1)$; these $n+1$ ordered pairs are
distinct. The other two upper bounds use bijections. The lower bound
already holds for all finite tripods, hence for this subclass as well.

These conclusions concern the specified fixed-degree verbose pullback
construction. Ordinary positive-degree persistence barcodes of these
ultrametric spaces are empty, so they give no counterexample to the usual
bottleneck metric. The established degree-zero result is also unaffected.

\section{Verification and research status}
The proof is exact and self-contained; it does not depend on numerical
optimization or a finite search over pullbacks. Two standard-library
Python implementations accompany this note.

\begin{center}
\begin{tabular}{@{}lrr@{}}\toprule
Implementation & Supporting checks & Matrix computations\\\midrule
\texttt{verify\_boundary.py} & 6,653 & 720 persistence reductions\\
\texttt{verify\_ranks.py} & 3,650 & 356 boundary ranks\\\bottomrule
\end{tabular}
\end{center}

The first program implements oriented boundary-column reduction over
$\mathbb Q,\mathbb F_2,\mathbb F_3,\mathbb F_5,\mathbb F_7$, retaining
all zero-length pairs. It checks the family for $1\le k\le6$, including
two tie orders. It also tests every positive four-fiber size vector with
$4\le |S|\le8$ in three fields, resulting in 5,226 pairwise comparisons.
The upper bounds and endpoint obstruction are reproduced exactly.
Its matching routine introduces no artificial diagonal points.

The second program does not import the first. It builds the boundary
matrices at thresholds $0,1,2$, uses dense row elimination over
$\mathbb Q$ and $\mathbb F_3$, checks vanishing positive-degree homology,
and compares rank increments with the claimed barcode multiplicities.
It additionally checks 3,267 finite instances of
\eqref{eq:all-multiplicities}. These finite tests do not establish the
universal quantifier; the subset-counting proof does.

Run, from the extracted package directory:
\begin{verbatim}
python3 verify_boundary.py --output boundary_results.json
python3 verify_ranks.py --output rank_results.json
\end{verbatim}

\textbf{Limitations.} Separate computational implementations are not
independent mathematical review. No Lean or other proof-assistant
certificate is supplied. The available repository and literature audit
found no matching prior 544 submission, but it is not an exhaustive
certification of unpublished work or inaccessible records. The accompanying resolution report is prepared for submission by
Siavash Sadeghi. Under the repository's evidence rules~\cite{contributing},
the appropriate status is \emph{Solution claimed} until a documented
independent audit or publication establishes a stronger status.

\begin{thebibliography}{9}
\bibitem{aim}
S. Brunton, M. Colbrook, M. de Hoop, G. Stepaniants, A. Townsend, and R. Ward,
\emph{AIM: Original problem 544, Triangle inequality for the pullback
distance of verbose persistence barcodes}. Target at
\href{https://github.com/MColbrook/AIM/blob/aa776a01d7d48a79f93251af11fde9454b0aea95/problems/544-verbose-persistence-pullback-triangle.md}{commit \texttt{aa776a0}};
\href{https://github.com/MColbrook/AIM/blob/fa98b7525fa3f78317536a8825f9cfa0ae1c369c/problems/534-verbose-persistence-pullback-triangle.md}{current AIM 534 statement}
inspected 2 October 2026.

\bibitem{mz}
F. M\'emoli and L. Zhou,
\emph{Ephemeral persistence features and the stability of filtered chain complexes},
Journal of Computational Geometry \textbf{15}(2), 258--328
(volume labeled 2024; published 2025).
\href{https://doi.org/10.20382/jocg.v15i2a8}{doi:10.20382/jocg.v15i2a8}.
See Remark~1.1; Definitions~4.7, 5.3; Example~5.10 and Figure~7;
and Proposition~6.2. Author revision:
\href{https://arxiv.org/abs/2208.11770v8}{arXiv:2208.11770v8}, 21 August 2025.

\bibitem{contributing}
\emph{AIM contribution and evidence rules}.
\href{https://github.com/MColbrook/AIM/blob/fa98b7525fa3f78317536a8825f9cfa0ae1c369c/CONTRIBUTING.md}{\texttt{CONTRIBUTING.md}},
inspected 2 October 2026.
\end{thebibliography}
\end{document}
